\subsection{Summary of the Physical Formulation}

This chapter has systematically formulated the rigorous physical and operational constraints governing individual Microgrids. Rather than adopting simplistic linear models, the proposed framework leverages the exact Branch Flow Model (DistFlow) to accurately capture the non-linear active and reactive power (P-Q-V) dynamics, preserving the strict operational bounds of local assets including dispatchable generators, dynamic BESS state-of-charge, and critical load protection via VOLL penalties.

To overcome the inherent non-convexity of the precise AC-OPF problem, the Second-Order Cone Programming (SOCP) relaxation was introduced. This mathematical convexification successfully transforms the intractable physical grid constraints into a convex inequality cone, guaranteeing rapid computational convergence while maintaining exact physical feasibility. Crucially, the model incorporates structural mechanisms to prevent exactness loss during emergency energy surplus conditions by mathematically prioritizing zero-cost renewable curtailment over fictitious line losses. By establishing a strictly convex and exact local physical formulation, this chapter fulfills the critical prerequisite for distributed market coordination. This convexified AC-OPF foundation fundamentally enables the application of the Analytical Target Cascading (ATC) algorithm, which will be developed in the following chapter to securely orchestrate P2P energy trading across the Multi-Microgrid system without compromising non-linear grid stability.
